\begin{theorem}[CauchySubsequenceExtractor regular-Cauchy forward link]
\label{thm:cauchy-subsequence-extractor-regular-cauchy-forward-link}
Every accepted $\CauchySubsequenceExtractorUp$ packet consumes the $\RegularCauchyMinModulusUp$ witness of \autoref{ch:concrete-instances-regular-cauchy-min-modulus-namecert}, the $\RegularCauchyCofinalSubsequenceUp$ cofinal route of \autoref{ch:concrete-instances-regular-cauchy-cofinal-subsequence-namecert}, and the selector dependency rows of \autoref{ch:concrete-instances-cauchy-subsequence-selector-namecert} before exporting a subsequence witness. Thus the extraction row $E$ is a finite regular-Cauchy forward link over the displayed window $W$ and dyadic ledger $D$, not an independent subsequence selector.
\end{theorem}
\begin{proof}
Project the carrier $S=(C,Q,W,D,E,R,L,H,T,P,N)$ from \autoref{def:cauchy-subsequence-extractor-carrier}. The extraction row $E$ is downstream of the Cauchy source $C$, regular readback $Q$, finite window $W$, and dyadic tolerance ledger $D$.

Use the selector-lattice handoff of \autoref{thm:cauchy-subsequence-extractor-selector-lattice-handoff}. It binds $E$ to the displayed selector schedule before any $\RealUp$ or $\RealSeparabilityUp$ consumer can read the terminal rows.

Now attach the regular-Cauchy routes. The minimum-modulus forward link forces the same $W,D,E$ rows to consume the $\RegularCauchyMinModulusUp$ witness, while the cofinal-subsequence dependency route sends the same extraction row to the $\RegularCauchyCofinalSubsequenceUp$ consumer before the inherited real seal is exposed.

The structural rows $H,T,P,N$ transport, replay, source, and name that single route. Since no carrier coordinate stores a free selector, countable choice, quotient stream equality, or ambient completeness, the exported subsequence witness is exactly the displayed regular-Cauchy forward link.
\end{proof}
